\documentclass[11pt]{article}
% Internal note (not part of any submitted paper).
\usepackage[letterpaper,margin=25mm]{geometry}
\usepackage{amsmath,amssymb,amsthm,bm}
\usepackage{booktabs}
\usepackage[round,authoryear]{natbib}
\usepackage{url}
\usepackage[hidelinks]{hyperref}
\newcommand{\E}{\mathbb E}

\title{Alignment-free bigram and trigram distances\\under the MixFrag Pair HMM}
\author{Internal note, tkf-dp}
\date{September 2026}

\begin{document}
\maketitle

\begin{abstract}
We derive closed-form expectations, under the MixFrag Pair HMM with a profile
mixture substitution model, of the number of shared bigrams and trigrams
between two unaligned sequences, and invert them to estimate evolutionary
time. The expectations match simulation to within about 2\%, and on
simulated pairs the trigram estimator is precise up to about 1.5
substitutions per site. On Pfam seed alignments, however, the estimators are
biased low relative to maximum-likelihood distances, and guide trees built
from them, including with delta-method variances in BioNJ, agree poorly with
reference trees. We record the derivation and the negative result.
\end{abstract}

\section{Motivation}

Progressive aligners need a guide tree before any alignment exists. MUSCLE
builds its first tree from the fraction of shared $k$-mers
\citep{Edgar2004}. For a statistical aligner it is natural to ask for an
alignment-free distance whose expectation is computed under the aligner's own
evolutionary model, so that it estimates time $t$ on the model's scale. The
model here has MixFrag indels \citep{HolmesMixFrag}, the TKF92
\citep{ThorneKishinoFelsenstein1992} fragment model with a mixture of fragment
length distributions, and a $K$-class profile mixture for substitutions.

\section{Model}

\paragraph{Substitution.} Each column has a latent class $k$ with weight
$w_k$; given $k$, residues evolve under the reversible chain
$Q^k_{xy}=r_kS_{xy}\pi^k_y$ ($x\ne y$), with exchangeabilities $S$ shared by
all classes, as in LG \citep{LeGascuel2008} combined with the C20 profiles
\citep{LeGascuelLartillot2008}. Write $M^k(t)=e^{tQ^k}$ and
$\bar\pi=\sum_kw_k\pi^k$.

\paragraph{Indels.} MixFrag is TKF92 in which each fragment draws a fragtype
$f$ with weight $w_f$ and then has geometric length with extension
probability $r_f$; the indel and substitution processes are shared across
fragtypes, and one fragtype gives TKF92. Fragments are born and die as units
under the TKF91 links process \citep{ThorneKishinoFelsenstein1991} with
insertion and deletion rates $\lambda<\mu$. Over time $t$, each ancestral
fragment survives with probability $\alpha=e^{-\mu t}$, and the number of
fragments inserted after a surviving one is positive with probability
\[
\beta=\frac{\lambda e^{-\lambda t}-\lambda\alpha}{\mu e^{-\lambda t}-\lambda\alpha}.
\]
The pair HMM is reversible, so for two present-day sequences either may be
treated as the ancestor of the other at distance $t$.

\section{Expected shared bigrams and trigrams}

\paragraph{Statistics.} For unaligned sequences $X$ and $Y$ with $j$-gram
counts $c^{(j)}_X$ and $c^{(j)}_Y$ (here $j=2$ or $3$), let
\[
D_j=\sum_wc^{(j)}_X(w)\,c^{(j)}_Y(w),
\]
the number of position pairs at which the $j$-grams starting there are
equal. Given the
alignment, residues in different columns are independent, so each letter of a
$j$-gram matches with probability
\[
\begin{cases}
s(t)=\sum_kw_k\sum_a\pi^k_aM^k_{aa}(t)&\text{if the two positions are homologous,}\\
\rho=\sum_a\bar\pi_a^2&\text{otherwise,}
\end{cases}
\]
where $\bar\pi=\sum_kw_k\pi^k$. $\E[D_j]$ is therefore a sum, over the
patterns of homology along $j$ consecutive diagonal position pairs, of the
expected number of occurrences of each pattern times
$s^{\#\text{homologous}}\rho^{\#\text{not}}$.

\paragraph{The Pair HMM supplies the patterns.} Take $X$ as the ancestor.
Each ancestral residue survives with probability $\alpha=e^{-\mu t}$.
Fragments are indivisible and insertions occur only between them, so two
adjacent ancestral residues in one fragment survive together with nothing
inserted between them; across a fragment boundary, a surviving residue is
followed directly by its surviving successor with probability
$(1-\beta)\alpha$ (no insertion, then survival). Following the fragtype of the
current residue along the ancestor, the probability that $j-1$ consecutive
adjacencies are all preserved, given that the first residue survives, is
\[
p_j=\omega^\top\mathsf T^{\,j-1}\mathbf 1,\qquad
\mathsf T_{fg}=r_f\,\delta_{fg}+(1-r_f)(1-\beta)\alpha\,w_g,\qquad
\omega_f\propto\frac{w_f}{1-r_f},
\]
where $\omega$ is the fragtype distribution of a residue in a long stationary
sequence. In particular $p_2=1-\big(1-(1-\beta)\alpha\big)/\bar\ell$, with
$\bar\ell=\sum_fw_f/(1-r_f)$ the mean fragment length. With
$n=\alpha(L_X-j+1)$ and $N_j=(L_X-j+1)(L_Y-j+1)$,
\begin{align}
\E[D_2]&\approx\rho^2N_2+n(s-\rho)\big(2\rho+(s-\rho)\,p_2\big),\label{eq:d2}\\
\E[D_3]&\approx\rho^3N_3+n\Big(p_3(s^3-\rho^3)+2(p_2-p_3)(s^2\rho-\rho^3)\notag\\
 &\hphantom{{}\approx\rho^3N_3+n\Big(}+(3-4p_2+p_3)(s\rho^2-\rho^3)\Big).\label{eq:d3}
\end{align}
The trigram expression neglects the rare pattern in which the outer two
position pairs are homologous but the middle one is not (a deletion and an
insertion side by side). The other approximations are the long-sequence limit
for $\omega$ and conditioning on the ancestor's length only; by reversibility
either sequence may play the ancestor, and we use the mean of the two lengths.

\paragraph{Estimators.} $\alpha$, $s-\rho$ and the $p_j$ all decrease in $t$,
so each expectation is monotone in $t$, and $\hat t$ is its root, found by
bisection. The cost per pair is $O(L_X+L_Y+A^j)$. Three remarks:
\begin{itemize}
\item Bigrams see MixFrag only through the mean fragment length $\bar\ell$,
  so a TKF92 model with the same $\bar\ell$ gives the same bigram distance.
  Trigrams, through $p_3$, see the spread of fragment lengths.
\item Real sequences depart from the model's composition, which inflates
  $D_j$. Replacing $\rho$ by the pair's own composition,
  $\rho_{\mathrm{obs}}=\sum_af_X(a)f_Y(a)$, corrects for this. The unigram
  analogue, $\E[D_1]-\rho L_XL_Y=\alpha L_X(s-\rho)$, cannot use this
  correction, because $D_1=\rho_{\mathrm{obs}}L_XL_Y$ exactly.
\item Longer $j$-grams have a larger signal per homologous run but a smaller
  background, and in simulation trigrams are the more precise of the two
  wherever either is informative.
\end{itemize}

\paragraph{Simulation.} We simulated sequence pairs by walking the MixFrag
Pair HMM ($\mu=0.05$, about 130 fragments or 300 residues,
$r=(0.25,0.75)$, $w=(0.65,0.35)$) with LG substitution, 300 pairs per $t$.
Equations \eqref{eq:d2} and \eqref{eq:d3} match the simulated means to within
0.5\% and 1\% at every $t$, and with the C20 mixture to within 0.7\% and
2.3\% (the largest deviation at $t=3$). Median $\hat t$
[interquartile range]:

\begin{center}
\footnotesize
\setlength{\tabcolsep}{3pt}
\begin{tabular}{lccccccc}
\toprule
$t$ & 0.1 & 0.25 & 0.5 & 1.0 & 1.5 & 2.0 & 3.0\\
\midrule
bigram  & 0.11 [.06,.16] & 0.27 [.20,.33] & 0.51 [.40,.66] & 0.97 [.77,1.4]
        & 1.5 [1.1,3.4] & 2.5 [1.4,20] & 7.4 [1.6,20]\\
trigram & 0.10 [.08,.12] & 0.25 [.22,.28] & 0.51 [.45,.56] & 1.02 [.88,1.2]
        & 1.6 [1.3,2.5] & 2.6 [1.6,20] & 20 [1.9,20]\\
\bottomrule
\end{tabular}
\end{center}

Both are usable to $t\approx1.5$; the trigram estimator has roughly half the
spread. Beyond $t\approx2$ the excess over the background is too small and
estimates spread to the search bound (20). The unigram estimator is much
noisier throughout.

\paragraph{Pfam.} We compared $\hat t$ with the maximum-likelihood $t$ under TKF92+LG-C20, given
each pair's alignment from the Pfam seed MSA, using the Pfam-trained TKF92
indel parameters ($\lambda=0.0311$, $\mu=0.0317$, $r=0.674$), on 1200 pairs
from 60 random families of the test split of the Pfam seed alignments. These pairs are
divergent: the maximum-likelihood $t$ has median 2.9 (IQR 1.5--5.0).

\begin{center}
\small
\begin{tabular}{lcccccc}
\toprule
& & \multicolumn{4}{c}{median $\hat t/t_{\mathrm{ML}}$, by $t_{\mathrm{ML}}$} & \\
\cmidrule(lr){3-6}
estimator & Spearman & $<1$ & 1--2 & 2--4 & 4--10 & at bound\\
\midrule
trigram, $\rho_{\mathrm{obs}}$ null & 0.86 & 0.76 & 0.61 & 0.46 & 0.60 & 15\%\\
bigram, $\rho_{\mathrm{obs}}$ null  & 0.81 & 0.99 & 0.75 & 0.62 & 1.46 & 19\%\\
trigram, model null                 & 0.85 & 0.73 & 0.58 & 0.43 & 0.57 & 14\%\\
bigram, model null                  & 0.69 & 0.75 & 0.53 & 0.40 & 0.62 & 16\%\\
unigram, model null                 & 0.32 & 0.39 & 0.29 & 0.20 & 0.57 & 27\%\\
\bottomrule
\end{tabular}
\end{center}

Trigrams rank pairs best (Spearman 0.86) and depend little on the choice of
null, but they underestimate $t$ throughout. Bigrams with the
composition-corrected null are unbiased for $t<1$ (IQR of the ratio
0.80--1.29) but underestimate by 25--40\% for $1<t<4$. Both saturate beyond
$t\approx4$. So real pairs share more $j$-grams than the model predicts, and
more so the longer the $j$-gram; a plausible cause, not yet tested, is that
conserved sites cluster along the sequence, which an independent-sites model
cannot express.

\section{Guide trees}

The simulations suggest that the trigram distance could order merges well.
We tested this on the same Pfam families (53 with at least 8 sequences, each
subsampled to at most 40), computing all pairwise distances within each
family.

\paragraph{Variance-weighted BioNJ.} BioNJ \citep{Gascuel1997} carries a
variance for every distance and, at each agglomeration, weights the two
merged nodes to minimize the variance of the new distances; its default
assumes variances proportional to distances. Our estimator provides its own
variances. By the delta method,
\[
\mathrm{Var}(\hat t)\approx\frac{\mathrm{Var}(D_3)}{\big(\partial_t\E[D_3]\big)^2},
\]
where the derivative is taken from \eqref{eq:d3} and $\mathrm{Var}(D_3)$ is
estimated by a moving-block bootstrap over the trigram positions of each
sequence (50 replicates, blocks of 10), which captures overdispersion from
repeats and composition. Estimates at the search bound are censored and
given effectively infinite variance. The variance grows roughly
exponentially with $t$, which is also the regime targeted by weighted
neighbor joining \citep{Bruno2000}.

\paragraph{Test.} The reference tree for each family is BioNJ on the pairwise
maximum-likelihood distances under TKF92+LG-C20, given the pairwise
alignments induced by the seed MSA. As a control we also built BioNJ trees on
Kimura distances computed from the seed MSA. Trees were compared by the
normalized Robinson--Foulds distance \citep{RobinsonFoulds1981}.

\begin{center}
\small
\begin{tabular}{lcccc}
\toprule
& \multicolumn{4}{c}{mean normalized RF to the reference}\\
\cmidrule(lr){2-5}
tree & all (53) & median $t_{\mathrm{ML}}<3$ (24) & 3--5 (17) & $>5$ (12)\\
\midrule
NJ, trigram distances              & 0.82 & 0.71 & 0.88 & 0.94\\
BioNJ, trigram distances           & 0.78 & 0.67 & 0.81 & 0.94\\
BioNJ, trigram distances and variances & 0.77 & 0.68 & 0.79 & 0.93\\
Control: BioNJ, Kimura distances from MSA & 0.49 & 0.45 & 0.51 & 0.56\\
\bottomrule
\end{tabular}
\end{center}

BioNJ improves on NJ (better in 26 families, worse in 10; sign test
$p=0.011$), and the delta-method variances improve on NJ (31 better, 12
worse; $p=0.005$) but not on BioNJ's default variances (21 better, 17 worse;
$p=0.63$). All three trigram trees are far from the reference, and far worse
than the alignment-based control, even in the least divergent families,
where only 5\% of trigram distances are censored.

\section{Conclusion}

The closed forms are accurate and the estimators behave as intended on data
simulated from the model, but not on real protein families: real pairs share
more $k$-grams than the model predicts, increasingly so for longer
$k$-grams, and the resulting trees are poor. A plausible cause, not tested
here, is that conserved sites cluster along real sequences, which an
independent-sites model cannot express. The Pfam seed families are also more
divergent (median maximum-likelihood $t$ about 3) than the range in which
these statistics are informative. For guide trees, alignment-based distances
remain the better choice.

\begin{thebibliography}{9}
\bibitem[Bruno et~al.(2000)]{Bruno2000}
W.~J. Bruno, N.~D. Socci and A.~L. Halpern. Weighted neighbor joining: a
likelihood-based approach to distance-based phylogeny reconstruction.
\emph{Mol.\ Biol.\ Evol.} 17(1):189--197, 2000.
\bibitem[Edgar(2004)]{Edgar2004}
R.~C. Edgar. MUSCLE: multiple sequence alignment with high accuracy and high
throughput. \emph{Nucleic Acids Res.} 32(5):1792--1797, 2004.
\bibitem[Gascuel(1997)]{Gascuel1997}
O.~Gascuel. BIONJ: an improved version of the NJ algorithm based on a simple
model of sequence data. \emph{Mol.\ Biol.\ Evol.} 14(7):685--695, 1997.
\bibitem[TKF-DP supplement(2026)]{HolmesMixFrag}
TKF-DP supplement, Appendix~A.3, ``The TKF92 model with fragment mixtures
(MixFrag)''. \url{https://tkfdp.net}.
\bibitem[Le and Gascuel(2008)]{LeGascuel2008}
S.~Q. Le and O.~Gascuel. An improved general amino acid replacement matrix.
\emph{Mol.\ Biol.\ Evol.} 25(7):1307--1320, 2008.
\bibitem[Le et~al.(2008)]{LeGascuelLartillot2008}
S.~Q. Le, O.~Gascuel and N.~Lartillot. Empirical profile mixture models for
phylogenetic reconstruction. \emph{Bioinformatics} 24(20):2317--2323, 2008.
\bibitem[Robinson and Foulds(1981)]{RobinsonFoulds1981}
D.~F. Robinson and L.~R. Foulds. Comparison of phylogenetic trees.
\emph{Math.\ Biosci.} 53:131--147, 1981.
\bibitem[Thorne et~al.(1991)]{ThorneKishinoFelsenstein1991}
J.~L. Thorne, H.~Kishino and J.~Felsenstein. An evolutionary model for
maximum likelihood alignment of DNA sequences. \emph{J.\ Mol.\ Evol.}
33(2):114--124, 1991.
\bibitem[Thorne et~al.(1992)]{ThorneKishinoFelsenstein1992}
J.~L. Thorne, H.~Kishino and J.~Felsenstein. Inching toward reality: an
improved likelihood model of sequence evolution. \emph{J.\ Mol.\ Evol.}
34(1):3--16, 1992.
\end{thebibliography}

\end{document}
